\chapter{Processes and Synchronisation}
\label{ch:processes}

\standfirst{The 1983 machinery---\ct{Process}, \ct{Semaphore}, \ct{Delay},
priorities---all still here, all still working, and now genuinely
concurrent.}

\section{Forking}

A block becomes a process by being sent \ct{fork}:

\begin{st}
[ self doSomethingSlow ] fork.
[ self doSomethingSlow ] forkAt: Processor userBackgroundPriority.
\end{st}

\ct{fork} answers nothing useful and \emph{swallows anything the block
raises}, which is tolerable for a background chore and wrong for a share of a
computation somebody is waiting on. For that, use \ct{forkParallel:}
(Chapter~\ref{ch:parallel}), which answers a \ct{Promise}.

To make a process without starting it:

\begin{st}
| p |
p := [ self work ] newProcess.
p priority: Processor userSchedulingPriority.
p resume.
\end{st}

\section{Priorities}

Eight of them. Seven have an accessor on \ct{Processor}:

\begin{center}
\small
\begin{tabular}{@{}lp{0.70\textwidth}@{}}
\toprule
8 & \ctp{timingPriority} --- the delay timer\\
7 & \ctp{highIOPriority}\\
6 & \ctp{lowIOPriority} --- the input process\\
5 & \ctp{userInterruptPriority}\\
4 & \ctp{userSchedulingPriority} --- the default, what your code runs at\\
3 & \ctp{userBackgroundPriority}\\
2 & \ctp{systemBackgroundPriority}\\
1 & the idle process. \ctp{SystemRockBottomPriority} is a class variable
  with no accessor, so unlike the seven above it there is no message
  that answers it\\
\bottomrule
\end{tabular}
\end{center}

A higher-priority process is \emph{preferred} when a worker picks work.
It is not exclusive, and it never was a lock here.

\begin{st}
Processor activePriority             "4, normally"
Processor activeProcess              "the running Process object"
Processor yield                      "let something else have a turn"
Processor activeWorkerIndex          "which native worker you are on"
Processor workerCount                "how many there are"
\end{st}

The last two are this system's, not 1983's. They are useful for partitioning
work and for making a test prove that it really ran in parallel. \ct{yield} is
a primitive here: it parks the process at the end of its ready list and lets
whatever is ready run. 1983's forked a helper process on every call, which was
fine on one thread and not on thirty-one.

\section{Process protocol}

\begin{center}
\small
\begin{tabular}{@{}ll@{}}
\toprule
\ctp{resume} & make it runnable\\
\ctp{suspend} & take it off the run queue\\
\ctp{terminate} & stop it for good\\
\ctp{priority} \ctp{priority:} & \\
\ctp{signalException:} & raise an exception inside it\\
\bottomrule
\end{tabular}
\end{center}

\section{Semaphores}

The Blue Book's primitive, and now genuinely atomic across native threads.

\begin{st}
| done |
done := Semaphore new.
[ self longJob. done signal ] fork.
done wait.                    "blocks until the signal arrives"
\end{st}

A semaphore counts. \ct{signal} increments, \ct{wait} decrements and blocks if
the count is zero. \ct{Semaphore forMutualExclusion} starts at one, so the
first \ct{wait} succeeds and the second blocks:

\begin{st}
| lock |
lock := Semaphore forMutualExclusion.
lock critical: [ "one process at a time in here" ]
\end{st}

\begin{stnote}[The bug that atomicity fixed]
The 1983 \ct{wait} tested \ct{excessSignals} and \emph{then} added itself to
the waiting list. A \ct{signal} landing between the two spent itself on an
empty list, and the waiter queued forever. On one thread that window does not
exist; on eight it is microseconds wide and perfectly reliable. Test-then-add
is now one critical section.
\end{stnote}

\ct{Semaphore} is the right tool for one-shot signalling and for counting a
resource. For mutual exclusion, prefer \ct{Mutex}: it releases in an
\ct{ensure:} and it detects re-entry instead of deadlocking silently.

\section{Delays}

\begin{st}
(Delay forSeconds: 2) wait.
(Delay forMilliseconds: 250) wait.
\end{st}

A delay expires while nothing is running, so it cannot be implemented by
polling in the idle loop---with every process waiting on one, every worker is
idle and a loop that only looks at what is already ready would look forever.
There is a dedicated VM timer thread holding the clock, and it posts to the
async signal queue when a deadline passes. The process that resumes delays
treats that signal as \emph{look at the clock}: it resumes whatever is due
and re-arms the timer for whatever is not, so a signal that arrives late, or
twice, is harmless. 1983's assumed every signal was the one it had armed for,
and two workers waiting on one-millisecond delays hung it every time.

\begin{stnote}
Two bugs were found building that, both worth knowing as shapes. A timer that
has fired but not yet delivered is \emph{still pending}: clearing the armed
flag before posting the signal leaves a window in which a waiter asks "is
anything still coming?", hears no, and gives up just before it arrives. And a
new producer turns latent races into real ones---the moment the timer existed,
ThreadSanitizer found two races in the async queue that were older than the
timer and had never been executed by any test before.
\end{stnote}

\section{Timing things}

\begin{st}
Time millisecondsToRun: [ self theWork ]
Time millisecondClockValue
Time now
Date today
\end{st}

\begin{stwarn}
\ct{Time now} and \ct{Date today} are 1983's in the \ct{st2026} profile. The
\ct{pharo-time} profile supersedes them with Pharo's Chronology package,
which is richer---\ct{3 seconds}, \ct{1 day}, \ct{DateAndTime now}---and
passes its own 633 tests. Which you have depends on the profile you built.
\end{stwarn}

\section{What a process actually is}

A \ct{Process} holds a chain of \ct{MethodContext} objects---stack frames, as
ordinary heap objects with a class and instance variables. Suspending one is
storing a pointer; there is no machine stack to unwind and no size to choose
in advance.

That is why processes here are cheap enough to make thousands of, and it is
also why the debugger can show you the stack, let you edit a method in the
middle of it, and restart a frame: the stack is data.

The scheduler multiplexes those green processes M:N over the native worker
pool. Deciding \emph{where} a process runs is the scheduler's business;
\ct{forkParallel:} says that work may proceed independently, not which thread
takes it.
